
\section{Conclusiones}

\subsection{Regulador Lineal LDO}

El regulador lineal diseñado cumple con las especificaciones de la tabla \ref{tabla:especificaciones_ldo}, sin embargo, no es ideal:

\begin{itemize}
	\item El rechazo de ripple de la fuente de alimentación no es bueno, y prácticamente permite el paso de toda la oscilación de entrada. Este punto debe ser analizado con mayor detalle al realizar la compensación del circuito.
	\item El lazo de corriente es lento, ya que cuenta con muchas etapas.
\end{itemize}

Asimismo, cuenta con algunas virtudes:
\begin{itemize}
	\item La regulación de línea es excelente, así como la regulación de carga.
	\item La eficiencia es muy buena para tratarse de un regulador lineal, sobre todo para resistencias de carga bajas.
	\item El valor de $V_{DO}$ es muy bueno.
\end{itemize}

Se concluye que el regulador es apto para la aplicación en este proyecto, pero se deberán tener precauciones durante el diseño de la etapa switching, o en su defecto se deberá robustecer el diseño.

Algunas reflexiones adicionales:
\begin{itemize}
    \item El limitador de corriente shunt actúa mas "rápido" que el limitador por foldback hecho con amplificadores operacionales, ya que estos últimos suelen ser "lentos".Esto permite que ante un corto (RL muy baja), no se dispare una elevada corriente y actúe primero el limitador de corriente.
    \item Un transistor de paso hecho con un PMOS presenta una tensión de Drop Out baja, este tipo de circuitos se los conoce como ''Low Drop Out''.
    \item Utilizar cargas activas en los amplificadores diferenciales presentan una mejor RRMC.
\end{itemize}

\subsection{Regulador BUCK}

El regulador funciona correctamente tras ciertas adaptaciones provisorias, que permiten anular una parte del ruido de conmutación que se filtra el circuito PWM.

Los resultados respetaron, en su mayoría, los requerimientos pedidos, a excepción de algunos valores particualres, como la eficiencia. Si bien se podría aspirar a valores cercanos al 95\%, no son malos para un primer prototipo.

La fuente podría usar algunas mejoras, como protecciones a la carga y a la alimentación, y para sí misma. Además, se requiere un estudio de propagación de ruido de conmutación, que fue lo que ocasionó graves problemas durante el desarrollo, e incluso tras las correcciones, sigue volviendo a la fuente inestable bajo ciertos escenarios.

Asimismo, si bien la compensación demuestra ser robusta en las simulaciones, se podría evaluar la posibilidad de aumentar el margen de fase a costa de ancho de banda, y de esta forma asegurar aún más la estabilidad.

Por otro lado, la elección del filtro \textit{Snubber} resultó ser insuficiente: no llega a amortiguar del todo los sobrepicos de tensión sobre el nodo de conmutación. Para corregir esto, bastará con elegir una capacidad más grande a costa de menor eficiencia (pero mayor robustez y estabilidad a la salida).

La eficiencia podrá ser mejorada fácilmente con el refinamiento del circuito de conmutación. En particular, la resistencia de \textit{pull-up} responsable de alimentar al gate driver, consume casi \SI{25}{mA} en estado bajo. También se puede jugar con la corriente de alimentación de los reguladores lineas, entre otras cosas, acercándonos a la combinación de valores más eficiente.
